\documentclass[11pt,a4paper]{article}
\usepackage[T1]{fontenc}
\usepackage{lmodern}
\usepackage[margin=24mm]{geometry}
\usepackage{amsmath,amssymb}
\usepackage{microtype}
\usepackage{fancyhdr,lastpage}
\usepackage[hidelinks]{hyperref}
\hypersetup{pdftitle={Gyires's permanent inequality for circulant matrices of order five: overview},pdfauthor={Dev D. Goyal}}
\setlength{\parindent}{0pt}
\setlength{\parskip}{4pt}
\raggedbottom
\pagestyle{fancy}
\fancyhf{}
\fancyhead[L]{\footnotesize\scshape Gyires's inequality for circulants of order five}
\fancyhead[R]{\footnotesize\scshape Dev D. Goyal}
\fancyfoot[R]{\footnotesize\thepage\,/\,\pageref*{LastPage}}
\renewcommand{\headrulewidth}{0.3pt}
\fancypagestyle{firstpage}{\fancyhf{}\renewcommand{\headrulewidth}{0pt}\fancyfoot[R]{\footnotesize\thepage\,/\,\pageref*{LastPage}}}
\newcommand{\per}{\operatorname{per}}
\newcommand{\T}{\mathsf{T}}
\begin{document}
\thispagestyle{firstpage}
{\LARGE\bfseries Gyires's permanent inequality for circulant matrices of order five\par}
\vspace{4pt}
{\large Overview of the result, its proof and its verification\par}
\vspace{6pt}
Dev D. Goyal \enspace\textperiodcentered\enspace October 2026 \enspace\textperiodcentered\enspace
\href{https://github.com/dvgyl/gyires-circulant-5}{\texttt{github.com/dvgyl/gyires-circulant-5}}
\par\bigskip
The repository contains a proof that Gyires's permanent inequality holds for every doubly stochastic circulant matrix
of order five, with its equality case, together with a stronger inequality for every nonnegative circulant of order
five and a sharp quantitative margin. The proof is presented in a self-contained paper and formalized in Lean~4 with
Mathlib; the formal proof uses only Lean's standard axioms, and all statements were also tested with exact computations
written independently of the formal proof.

The result was obtained in an AI research effort directed by Dev~D.~Goyal. ChatGPT (OpenAI, GPT-6 with Pro
reasoning), working autonomously from Dev~D.~Goyal's prompts, proposed the problem, found the proof and wrote the Lean
formalization. Claude Opus~5.5, used with Claude Code, checked the work independently and wrote the paper in its
present form. Dev~D.~Goyal has reviewed the proofs; this review is not an independent expert validation. The paper has
not been peer reviewed.

\section*{1. The conjecture}
A real square matrix is doubly stochastic if its entries are nonnegative and its row and column sums are all $1$; the
set of such $n\times n$ matrices is $\Omega_n$, and $J_n$ is the matrix all of whose entries are $1/n$. B.~Gyires
conjectured in 1978 that every $A\in\Omega_n$ satisfies
\[
\frac{4\per(A)^2}{\per(AA^\T)+\per(A^\T A)+2\per(A^2)}\ \ge\ \frac{n!}{n^n},
\]
with equality only at $A=J_n$ (Conjecture~20 in the survey of Cheon and Wanless, Linear Algebra Appl.\ 403, 2005). The
case $n=2$ is elementary, the inequality for $n=3$ was proved by D.~K.~Chang in 1990, and the case $n=4$ was proved by
C.~Ozols in 2026 with a computer-assisted proof formalized in Lean. The conjecture is open for $n\ge5$.

\section*{2. The result}
A circulant matrix of order five has entries $c_{ij}=w_{j-i}$, indices modulo five. The doubly stochastic circulants are
the convex combinations of the five powers of the cyclic permutation matrix: a four-dimensional simplex inside the
sixteen-dimensional polytope $\Omega_5$; this simplex contains matrices with zero entries. For every circulant $C$ of
order five with nonnegative weights (the entries of its first row):
\begin{itemize}
\item the strong inequality $24\per(CC^\T)\le625\per(C)^2$ holds, with equality exactly when all five weights are
equal;
\item if $z_1,\dots,z_4$ are the gaps between the sorted weights, then
$625\per(C)^2-24\per(CC^\T)\ge484\,(z_1^{10}+z_2^{10}+z_3^{10}+z_4^{10})$, and $484$ is best possible;
\item Gyires's inequality holds for $C$, with equality exactly when all weights are equal. In particular, for every doubly
stochastic circulant of order five the ratio of Section~1 is at least $5!/5^5=24/625$, with equality exactly at
$J_5$.
\end{itemize}

\section*{3. The proof}
Every circulant is normal, and $\per(C^2)\le\per(CC^\T)$ by a comparison of the form
$2\per(MN^\T)\le\per(MM^\T)+\per(NN^\T)$, so Gyires's inequality follows from the strong inequality. The defect
$625\per(C)^2-24\per(CC^\T)$ is a homogeneous polynomial of degree ten in the weights, invariant under the twenty affine
changes of indices $i\mapsto ai+b$ modulo five. These reduce the $120$ orderings of the weights to six. For each of the six,
the defect, written in the smallest weight and the gaps between consecutive sorted weights, has $1001$ coefficients,
all nonnegative integers: $996$ are positive, the least being $484$, and five are zero. This gives the
inequality, the margin and the equality case at once.

\section*{4. Verification}
The Lean project builds from source with Lean~4.35.0 and Mathlib~v4.35.0. Nine principal theorems are restated with
Mathlib's definitions only. Comparator, a Lean tool that checks that a proof proves exactly a statement given in a
separate file, confirms that the proved theorems have exactly these statements, and the proofs are accepted by Lean's
kernel and by the independent NanoDa and con-ron kernels. The chamber certificate is checked in
Lean coefficient by coefficient, by the \texttt{ring} and \texttt{positivity} tactics, without evaluation by compiled
code. A separate PARI/GP script recomputes every permanent from the definition, recomputes the chamber polynomials for
all $120$ orderings, and tests every statement of the paper.

\section*{5. The paper}
\textbf{Gyires's permanent inequality for nonnegative circulant matrices of order five.}\par
Conjecture~20 of Cheon and Wanless for doubly stochastic circulants of order five, with its equality case; the strong
inequality and its sharp margin for nonnegative circulants.\enspace\emph{(Special case proved.)}\par
{\footnotesize\href{https://github.com/dvgyl/gyires-circulant-5/blob/main/preprints/Gyires-inequality-for-circulants-of-order-five/paper.pdf}{Paper}%
\enspace\textperiodcentered\enspace
\href{https://github.com/dvgyl/gyires-circulant-5/blob/main/lean/Challenges/GyiresCirculant5.lean}{Lean statements}%
\enspace\textperiodcentered\enspace
\href{https://github.com/dvgyl/gyires-circulant-5/blob/main/verification/README.md}{Checks}\par}

The result concerns circulant matrices of order five only; Gyires's conjecture for all of $\Omega_5$ remains open. The
searches for earlier work were limited, and no claim of novelty or priority is made.
\end{document}
